\documentclass{amsart}
\usepackage{amsmath,amssymb,amsthm,hyperref}
\newtheorem{proposition}{Proposition}
\newtheorem{lemma}{Lemma}
\theoremstyle{definition}
\newtheorem{remark}{Remark}
\begin{document}

\title{Generalizability and statistical power of a single-organization
team study ($n=51$): a formal assessment}
\maketitle

\section*{Setup and what is to be argued}

The object under study is an empirical investigation of $n=51$ work teams
drawn from a \emph{single} organization, reporting associations among three
team-level constructs: psychological safety $X$, team learning behavior $M$,
and team performance $Y$. The implicit substantive model is a mediation
chain
\begin{equation}\label{eq:model}
  X \;\longrightarrow\; M \;\longrightarrow\; Y,
\end{equation}
i.e.\ psychological safety promotes learning behavior, which in turn
improves performance, possibly with a residual direct effect $X\to Y$.

The task is not a single closed-form theorem but a structured, rigorous
assessment of (i) external-validity threats from single-organization
sampling, (ii) the inferential limits imposed by $n=51$, and (iii) which
multi-organization design features most improve validity. We make the
statistical components precise as small propositions with proofs, and we
ground the substantive (research-design) components in those results. The
methodological logic follows the standard Campbell--Stanley/Shadish--Cook--Campbell
framework for validity and classical large-sample estimation theory.

We adopt the convention that ``proof'' here means: (a) a precise statement of
each quantitative claim, (b) a deductively valid derivation from stated
modeling assumptions, and (c) numerical instantiation at $n=51$.

\section*{Part (i): External validity under single-organization sampling}

\subsection*{A sampling-frame formalization.}
Let $\mathcal{P}$ denote the target population of work teams to which one
wishes to generalize: teams across industries, national contexts, team
types, and measurement regimes. Let $\theta$ be a parameter of interest,
e.g.\ the standardized indirect effect $a b$ in \eqref{eq:model}, where $a$
is the $X\to M$ slope and $b$ the $M\to Y$ slope (controlling for $X$).
A study samples not from $\mathcal{P}$ but from a sub-population
$\mathcal{P}_0\subset\mathcal{P}$ consisting of teams in one organization,
characterized by a fixed bundle of contextual moderators
$c_0=(\text{industry},\ \text{nation},\ \text{team type},\ \text{HR/culture
regime},\ \text{measurement protocol})$.

\begin{proposition}[External-validity bias is a moderation phenomenon]
\label{prop:ev}
Suppose the population effect depends on context $c$ through a moderated
parameter $\theta(c)$, and the study consistently estimates
$\theta(c_0)$ (the effect at the single organization's context $c_0$).
Let the target estimand be the population-average effect
$\bar\theta=\mathbb{E}_{c\sim\Pi}[\theta(c)]$ under the distribution $\Pi$
of contexts over $\mathcal{P}$. Then the asymptotic generalization bias is
\[
  \operatorname{bias} \;=\; \theta(c_0)-\bar\theta
   \;=\; \big(\theta(c_0)-\theta(\bar c)\big)
        \;-\;\tfrac12\,\operatorname{tr}\!\big(\nabla^2\theta(\bar c)\,
        \Sigma_c\big)+o(\|\Sigma_c\|),
\]
where $\bar c=\mathbb{E}_\Pi[c]$ and $\Sigma_c=\operatorname{Cov}_\Pi(c)$.
In particular, the bias is zero for all $c_0$ \emph{only if} $\theta$ does
not vary with context (no moderation); otherwise it is nonzero and is
\emph{not} reduced by increasing the within-organization sample size $n$.
\end{proposition}

\begin{proof}
By definition $\operatorname{bias}=\theta(c_0)-\bar\theta$. Add and subtract
$\theta(\bar c)$: $\operatorname{bias}=(\theta(c_0)-\theta(\bar c)) +
(\theta(\bar c)-\bar\theta)$. A second-order Taylor expansion of
$\theta$ about $\bar c$ gives, for $c$ near $\bar c$,
$\theta(c)=\theta(\bar c)+\nabla\theta(\bar c)^{\!\top}(c-\bar c)
+\tfrac12 (c-\bar c)^{\!\top}\nabla^2\theta(\bar c)(c-\bar c)+r(c)$ with
$r(c)=o(\|c-\bar c\|^2)$. Taking $\mathbb{E}_\Pi$, the linear term vanishes
because $\mathbb{E}_\Pi[c-\bar c]=0$, leaving
$\bar\theta=\theta(\bar c)+\tfrac12\operatorname{tr}(\nabla^2\theta(\bar c)\Sigma_c)+o(\|\Sigma_c\|)$.
Substituting yields the displayed formula. If $\theta$ is constant in $c$,
both $\theta(c_0)-\theta(\bar c)=0$ and $\nabla^2\theta\equiv0$, so the bias
vanishes for every $c_0$. Conversely, if $\theta$ is non-constant there exist
contexts $c_0$ with $\theta(c_0)\neq\theta(\bar c)$, and since the entire
expression depends on $c_0$ only through the deterministic offset
$\theta(c_0)-\theta(\bar c)$ (independent of $n$), no increase in $n$ within
$\mathcal{P}_0$ removes it: $\widehat\theta\to\theta(c_0)$ a.s.\ as
$n\to\infty$ by consistency, so the limit error remains
$\theta(c_0)-\bar\theta$.
\end{proof}

\begin{remark}[Substantive reading]
Proposition~\ref{prop:ev} formalizes the central external-validity threat:
single-organization sampling fixes the moderator vector at $c_0$, so the
estimate is unbiased for the \emph{local} effect $\theta(c_0)$ but carries a
\emph{fixed} bias for the population-average effect $\bar\theta$ whenever the
psychological-safety$\to$learning$\to$performance relationships are moderated
by industry, national culture (e.g.\ power distance affecting
voice/risk-taking), team type (knowledge vs.\ production teams), or
measurement method. The decisive qualitative consequence is that this bias
is \emph{irreducible by sample size}: collecting more teams from the same
organization shrinks sampling variance but leaves the systematic
generalization error untouched. This is precisely why external validity is a
\emph{design} problem (sample composition), not a \emph{power} problem.
\end{remark}

\subsection*{Common-method and range-restriction threats.}
Two further single-organization threats are formalizable.

\emph{(a) Common-method variance.} If $X$, $M$, $Y$ are all measured by the
same self-report instrument from the same raters, an additive method factor
$f$ inflates observed correlations. Writing observed scores
$\tilde V = V + \lambda_V f$ with $f\perp V$ and $\operatorname{Var}(f)=\sigma_f^2$,
\[
  \operatorname{Cov}(\tilde X,\tilde Y)
   = \operatorname{Cov}(X,Y) + \lambda_X\lambda_Y\sigma_f^2,
\]
so the observed covariance is biased upward by $\lambda_X\lambda_Y\sigma_f^2\ge0$
when loadings share sign. Single-method, single-source designs thus
systematically overstate the associations of interest; this bias, like the
moderation bias, does not vanish as $n\to\infty$.

\emph{(b) Range restriction.} Teams in one organization share selection,
training, and culture, compressing the variance of $X$. For a bivariate
normal pair, if the unrestricted correlation is $\rho$ and the predictor SD
is reduced by a factor $u=\sigma_{\text{restricted}}/\sigma_{\text{full}}<1$,
the observed correlation is
\[
  \rho_{\text{obs}}=\frac{u\,\rho}{\sqrt{1-\rho^2(1-u^2)}}<\rho,
\]
the classical Thorndike Case~II formula, which attenuates toward $0$ as
$u\to0$. Hence single-organization homogeneity can also \emph{deflate}
estimated relationships, and the direction of the net single-organization
bias (method inflation vs.\ range deflation) is itself not identifiable from
within-organization data.

\section*{Part (ii): Inferential limits at $n=51$}

\subsection*{Precision of a single correlation.}
\begin{proposition}[Wide confidence intervals]\label{prop:ci}
For a bivariate-normal correlation estimated from $n$ teams, Fisher's
$z=\tfrac12\log\frac{1+r}{1-r}$ is approximately $N\!\big(\tfrac12\log\frac{1+\rho}{1-\rho},\,\frac1{n-3}\big)$.
The two-sided $95\%$ CI for $\rho$ therefore has $z$-scale half-width
$1.96/\sqrt{n-3}$. At $n=51$ this is $1.96/\sqrt{48}=0.283$, and back-transforming,
a point estimate $r=0.40$ yields the interval $(0.14,\,0.61)$ (width $\approx0.47$),
and $r=0.30$ yields $(0.03,\,0.53)$ (width $\approx0.50$).
\end{proposition}

\begin{proof}
The variance-stabilizing property $\operatorname{Var}(z)\approx 1/(n-3)$ is
standard (delta method on the Fisher transform). The half-width is
$z_{0.975}/\sqrt{n-3}=1.96/\sqrt{48}=0.2829$. For $r=0.40$,
$z=0.4236$, giving the $z$-interval $(0.1407,0.7065)$; inverting
$\rho=\tanh(z)$ gives $(0.140,0.609)$. For $r=0.30$, $z=0.3095$,
$z$-interval $(0.0266,0.5924)$, $\rho$-interval $(0.027,0.532)$. Widths
$0.469$ and $0.505$ respectively. (Numerically verified.)
\end{proof}

\subsection*{Power.}
\begin{proposition}[Low power for small/medium effects]\label{prop:power}
For the two-sided $\alpha=0.05$ test $H_0:\rho=0$ using Fisher's $z$, the
power at true correlation $\rho$ and sample size $n$ is
\[
  \beta_{\text{pow}}(\rho,n)
   =\Phi\!\Big(\zeta\sqrt{n-3}-z_{0.975}\Big)
   +\Phi\!\Big(-\zeta\sqrt{n-3}-z_{0.975}\Big),
   \quad \zeta=\tfrac12\log\tfrac{1+\rho}{1-\rho}.
\]
At $n=51$: $\beta_{\text{pow}}(0.2)=0.29$, $\beta_{\text{pow}}(0.3)=0.57$,
$\beta_{\text{pow}}(0.4)=0.84$, $\beta_{\text{pow}}(0.5)=0.97$.
Conversely, achieving power $0.80$ requires $n\ge 783$ for $\rho=0.1$,
$n\ge 194$ for $\rho=0.2$, and $n\ge 85$ for $\rho=0.3$.
\end{proposition}

\begin{proof}
Under the alternative, $z\sim N(\zeta,1/(n-3))$, so $\sqrt{n-3}\,z\sim
N(\sqrt{n-3}\,\zeta,1)$ and the rejection region is
$|\sqrt{n-3}\,z|>z_{0.975}$; the displayed power is the probability of this
event for a normal with mean $\sqrt{n-3}\,\zeta$, giving the two-tail sum.
The required-$n$ formula solves
$\sqrt{n-3}\,\zeta=z_{0.975}+z_{0.80}$, i.e.
$n=3+\big((z_{0.975}+z_{0.80})/\zeta\big)^2$, with
$z_{0.975}=1.96,\ z_{0.80}=0.8416$. Plugging in $\rho=0.1,0.2,0.3$
($\zeta=0.1003,0.2027,0.3095$) gives $n=783,194,85$. The power values at
$n=51$ follow by direct substitution. (Numerically verified.)
\end{proof}

\subsection*{Multivariate inference and degrees of freedom.}
\begin{proposition}[Degrees-of-freedom scarcity]\label{prop:df}
Consider estimating the mediation model \eqref{eq:model} with $p$
team-level covariates by OLS regressions. The residual degrees of freedom in
the outcome regression are $n-p-1$. With $n=51$:
\begin{itemize}
\item Estimating $b$ (effect of $M$ on $Y$) while controlling for $X$ and
$k$ further covariates uses $p=k+2$ predictors, leaving $49-k$ residual df;
each additional moderator/control costs one df and inflates
$\operatorname{Var}(\hat\beta_j)=\sigma^2\,[(\mathbf X^\top\mathbf X)^{-1}]_{jj}
=\dfrac{\sigma^2}{(n-1)s_j^2(1-R_j^2)}$, where $R_j^2$ is the multiple
$R^2$ of predictor $j$ on the others (the variance-inflation factor
$\mathrm{VIF}_j=1/(1-R_j^2)$).
\item Testing the indirect effect $ab$ requires either the (poorly
calibrated at small $n$) Sobel normal approximation or the bootstrap; the
bootstrap percentile interval for $ab$ inherits the wide-CI behavior of
Proposition~\ref{prop:ci} because its width is $O(1/\sqrt{n})$.
\end{itemize}
Consequently, with $n=51$ one cannot reliably (a) fit models with more than a
handful of covariates without severe variance inflation, (b) detect
moderation (interaction terms, which are typically lower-powered than main
effects), or (c) distinguish full from partial mediation, since the
direct-effect estimate $c'$ shares the wide CIs above.
\end{proposition}

\begin{proof}
The OLS variance formula and df accounting are standard. For the variance of
$\hat\beta_j$, partition $\mathbf X$ and use the Frisch--Waugh--Lovell
theorem: $\hat\beta_j$ equals the slope of $Y$ on the part of $X_j$ orthogonal
to the other predictors, whose sum of squares is $(n-1)s_j^2(1-R_j^2)$;
hence $\operatorname{Var}(\hat\beta_j)=\sigma^2/[(n-1)s_j^2(1-R_j^2)]$, which
diverges as $R_j^2\to1$ (collinearity, increasingly likely as $p$ grows
relative to $n$). The $O(1/\sqrt n)$ width of the bootstrap interval for the
smooth functional $ab$ follows from the bootstrap CLT for plug-in estimators
of differentiable functionals (Efron--Tibshirani). Interaction-term power is
governed by Proposition~\ref{prop:power} applied to the partial correlation
of the product term, whose effective magnitude is typically smaller, hence
lower power.
\end{proof}

\subsection*{Multiplicity.}
Reporting several bivariate and partial associations among $\{X,M,Y\}$ plus
controls without multiplicity correction inflates the family-wise error rate:
for $m$ independent tests at level $\alpha$, $\mathrm{FWER}=1-(1-\alpha)^m$,
e.g.\ $m=5\Rightarrow 0.23$. At $n=51$ the joint problem of low power and
uncorrected multiplicity is acute: true effects are likely missed while
spurious ones may be reported, and surviving estimates are subject to the
winner's-curse/magnitude inflation typical of low-powered designs.

\section*{Part (iii): Design features that most improve validity}

We now show \emph{why} the proposed multi-organization features address
exactly the deficiencies proved above, and rank them by the threat they
remediate.

\paragraph{(D1) Stratified multi-organization sampling (addresses
Proposition~\ref{prop:ev}).}
Sampling teams from organizations spanning a designed grid of contexts
$c$ replaces the point estimand $\theta(c_0)$ with the population-average
$\bar\theta$ and, crucially, makes the moderators \emph{estimable}.
Formally, with strata indexed by $c$ and stratum weights $w_c$ matching the
target $\Pi$, the stratified estimator
$\widehat{\bar\theta}=\sum_c w_c\,\widehat\theta(c)$ is consistent for
$\bar\theta$, eliminating the irreducible bias of
Proposition~\ref{prop:ev}. Moreover, modeling $\theta(c)=\theta_0+\gamma^\top c$
turns the previously \emph{unidentified} moderation into estimable
coefficients $\gamma$, converting an external-validity threat into a
measured effect-size of generalizability. This is the single most important
improvement, because it is the only one that removes a bias that sample size
alone cannot.

\paragraph{(D2) Replication across cultural/national contexts (a special
case of D1, addresses both bias and generalizability inference).}
Cross-cultural replication estimates $\theta(c)$ at distinct national
contexts, directly testing whether
$\nabla^2\theta$ and the offset $\theta(c_0)-\theta(\bar c)$ in
Proposition~\ref{prop:ev} are nonzero. It also provides independent estimates
whose agreement (or disagreement) bounds the moderation term empirically,
yielding heterogeneity statistics ($\tau^2$, $I^2$) in a random-effects
meta-analytic combination.

\paragraph{(D3) Increased sample size $N$ (addresses
Propositions~\ref{prop:ci}--\ref{prop:df}).}
CI half-widths and required-$N$ scale as $1/\sqrt{N-3}$. Increasing
from $51$ to $500$ teams shrinks every $z$-scale CI by the factor
$\sqrt{48/497}=0.31$ (a $\sim$69\% reduction); to $5000$ teams by the factor
$0.098$. Power for $\rho=0.2$ rises from $0.29$ at $n=51$ to $\approx0.99$ at
$n=500$. Larger $N$ also restores degrees of freedom
(Proposition~\ref{prop:df}), enabling multivariate models with several
controls, well-powered interaction tests, and reliable bootstrap mediation
intervals. Note the division of labor: \emph{$N$ buys precision and power;
it does not buy external validity}, which only D1/D2 provide.

\paragraph{(D4) Multi-method, multi-source measurement (addresses the
common-method and range-restriction biases of Part (i)).}
Measuring $X,M,Y$ from different sources/instruments (e.g.\ peer-rated safety,
behaviorally coded learning, objective performance) drives the shared method
loading $\lambda$ toward zero, removing the
$\lambda_X\lambda_Y\sigma_f^2$ inflation term. Multi-organization sampling
simultaneously widens the range of $X$, raising the restriction ratio
$u\to1$ and undoing the Thorndike attenuation. A multitrait--multimethod
design lets the method variance itself be estimated and partialed out.

\paragraph{(D5) Multilevel modeling and clustering (a correctness, not just
power, requirement at scale).}
With teams nested in organizations, observations within an organization are
correlated. The design effect $\mathrm{DEFF}=1+(\bar m-1)\rho_{\text{ICC}}$
(average cluster size $\bar m$, intraclass correlation $\rho_{\text{ICC}}$)
gives effective sample size $N_{\text{eff}}=N/\mathrm{DEFF}$. Ignoring
clustering understates standard errors; a hierarchical (mixed) model with
organization-level random effects both corrects inference and \emph{estimates}
the between-organization variance that quantifies generalizability. (At
$n=51$ from one organization this variance is, by construction,
unidentified -- a structural limitation, not a power limitation.)

\subsection*{Ranking.}
The remediation hierarchy follows from which deficiency each feature removes:
\begin{enumerate}
\item \textbf{D1/D2 (stratified, cross-context sampling)} -- uniquely remove
the \emph{irreducible} external-validity bias of
Proposition~\ref{prop:ev} and make generalizability itself estimable. Most
important.
\item \textbf{D4 (multi-method/source)} -- remove the
sample-size-invariant method-inflation and range-restriction biases of
Part (i).
\item \textbf{D3 (larger $N$)} -- remove the variance/power deficiencies of
Propositions~\ref{prop:ci}--\ref{prop:df}; necessary but \emph{not
sufficient} for generalization.
\item \textbf{D5 (multilevel modeling)} -- ensures the inference at scale is
\emph{correct} and turns between-organization heterogeneity into a measured
quantity.
\end{enumerate}

\section*{Conclusion}

We have established the following rigorous claims.
(i) Single-organization sampling produces an external-validity bias
$\theta(c_0)-\bar\theta$ that is nonzero precisely when the
safety--learning--performance relationships are context-moderated, and this
bias is \emph{invariant to within-organization sample size}
(Proposition~\ref{prop:ev}); common-method variance and range restriction
add further sample-size-invariant distortions of opposite sign whose net
effect is not identifiable from within-organization data.
(ii) At $n=51$, confidence intervals for the reported correlations are wide
($\approx0.47$--$0.50$ on the correlation scale; Proposition~\ref{prop:ci}),
power for small-to-medium effects is low ($0.29$ at $\rho=0.2$, $0.57$ at
$\rho=0.3$; Proposition~\ref{prop:power}), and degrees-of-freedom scarcity
sharply limits multivariate, moderation, and mediation inference
(Proposition~\ref{prop:df}), with uncontrolled multiplicity compounding the
problem.
(iii) Stratified multi-organization and cross-cultural sampling (D1/D2) are
the only design features that remove the irreducible external-validity bias
and render generalizability estimable; multi-method/multi-source measurement
(D4) removes the method and range biases; larger $N$ (D3) supplies the
precision, power, and degrees of freedom that $51$ teams lack; and
hierarchical modeling (D5) makes large-sample inference valid and quantifies
between-organization heterogeneity.

Hence the reported single-organization relationships are best read as
\emph{locally} estimated effects $\theta(c_0)$ with poor precision, whose
generalization beyond the studied context requires designed contextual
variation, not merely a larger homogeneous sample.

\end{document}
